\section{Discussion}
\label{sec:discussion}

\subsection{Whether one policy can be best across playstyles}

The answer has two parts, one for each half of the agent: the inference
component, derived from the rules, and the decision component, fitted against a
panel of opponents.

\textbf{The inference half is playstyle-independent.} Constraints (C1)--(C5) of
\S\ref{sec:belief} are logical consequences of the dialect specified in
\S\ref{sec:rules}, so they are true whatever policy the observed player is
running, and no opponent behaviour can make them false. Nothing in the
observer-conditioned deal posterior is fitted to an opponent, and the same is true
of Theorem~\ref{thm:locked}, which is a statement about which asks the rules
permit rather than about how anyone chooses among them. Consistent with this, the
frozen-policy ablations that remove parts of the constraint set produce the
largest win-rate decreases recorded in Table~\ref{tab:ablations}: dropping the
ask-legality certificates gives a full-minus-ablated difference of
\numAblSinkhorn{} points, with a 95\% paired percentile bootstrap interval over
deals of \numAblSinkhornCI{}. That row removes the fitted prior along with the
certificates, and it is a frozen-policy sensitivity rather than the standalone
causal value of either (\S\ref{sec:ablations}).

\textbf{The decision half cannot be a best response to every opponent at once.}
Choosing among asks trades off quantities --- the cost of handing a particular
opponent the turn, the information an ask reveals, when to declare a half-suit
--- whose values depend on what the other five seats do next. A best response to
a policy is a function of the policy it responds to, so no single deterministic
policy is simultaneously a best response to every opponent, and no experiment in
this paper could establish that one was. The defensible target is therefore
minimal exploitability rather than dominance over an arbitrary opponent. For a
game with two three-player teams, no communication channel, and a shared team
payoff, the appropriate solution concept is a team-maxmin equilibrium with
correlation \cite{celli-gatti,zhang-tbdag}, which we do not compute and which is
out of computational reach at this game size. The fitting objective
\eqref{eq:softmin} is a deliberate finite-panel approximation of it: a soft
minimum over a hand-built six-opponent panel rather than a minimum over the whole
policy space.

\textbf{What the evidence supports.} Within the eight-policy panel of
Table~\ref{tab:h2h}, one policy is ahead of every member simultaneously:
\vfast's lowest win rate against any panel member is \numWorstCase\%, recorded
against \vthree{} on the primary bank under the dialect of \S\ref{sec:rules} and
its default arbitration order, two alternatives to which put it slightly lower
(\S\ref{sec:robustness}).
That statement is bounded by the panel, and it measures generalisation to new
deals rather than to new strategic populations (\S\ref{sec:results-h2h}). The
local response probe is inconclusive in the other direction: a response fitted
specifically to frozen \vfast{} reached \numLBRFour\% with a 95\% deal-clustered
percentile bootstrap interval of \numLBRFourCI{} that includes 50\%, and an
achieved exploiter score is in any case a lower bound within the searched class
(\S\ref{sec:lbr}). The two results are not in tension: the fitted policy is
uniformly ahead of the opponent population we were able to construct, and its
exploitability by a purpose-built response remains unresolved.

\subsection{Interpreting the belief-substitution result}

Substituting the exact reference belief into a policy fitted around Fast beliefs
cost \numAblBlock{} points, with a 95\% paired percentile bootstrap interval over
deals of \numAblBlockCI{} (E5, \S\ref{sec:ablations}). The fitted policy is
therefore not invariant to the belief representation; that is not evidence that
an exact-belief policy would be weaker after matched retraining, and no matched
fit was run. Two mechanisms are consistent with the observation, and this study
cannot separate them.

The first is threshold placement. The \numKnobs{} decision knobs of
\S\ref{sec:policy} --- among them the declaration threshold, the minimum team
probability, the two cheap gate thresholds and the stopping margin --- were
fitted against the distribution of scores produced by the Fast estimator. E2
records a mean absolute difference of \numSinkhornMean{} and a maximum of
\numSinkhornMax{} between Fast and block per-card marginals, so the reference
engine presents a differently distributed input to knobs placed at quantiles of
the Fast distribution. Under that account the loss is a calibration mismatch
between fitted constants and a substituted estimator, and refitting would remove
it.

The second is rank tolerance. The ask rule scores the legal ask candidates
enumerated each turn and takes an argmax, which is insensitive to perturbations
that do not reorder the leaders, so a more accurate estimator can change the
chosen ask only where the ranking is close. Whether the largest Fast--Block
deviations coincide with close rankings is a conjecture rather than a
measurement: E2 records the size of the deviations but no position-conditioned
statistic. E12 does not resolve the question either --- under a deliberately
unfitted posterior-greedy asking rule the two belief paths reached
\numETwelveFast\% and \numETwelveBlock\%, with 95\% deal-clustered percentile
bootstrap intervals of \numETwelveFastCI{} and \numETwelveBlockCI{}, a gap in the
same direction as the fitted comparison but with substantially overlapping
intervals.

Settling the question needs the matched-budget pair of fits described in
\S\ref{sec:limitations}: the cross-entropy procedure of \S\ref{sec:fitting} run
twice from the same seeds and generation count, once with each inference path in
the inner loop. At the reference engine's \numThroughputBlock{} games per second
(E9) that is expensive but not infeasible; it was not attempted here.

\subsection{Implications for human play}
\label{sec:humanplay}

Three of the mechanisms this analysis relies on are cheap enough to apply at a
table without a posterior over deals. An ask is a confession: under the dialect
of \S\ref{sec:rules} it establishes that the asker lacks the named card and holds
some other card of its half-suit, and a miss establishes that the target lacks
the named card, which are constraints (C1)--(C3) and (C5) of \S\ref{sec:belief}
applied by hand. Public hand sizes then close those exclusions into certainties
--- the capacity constraint (C4) applied by hand --- and the single frozen-policy
ablation that removes capacity conditioning produced the largest win-rate
decrease in Table~\ref{tab:ablations}. A half-suit held entirely by one team
cannot change hands, so waiting carries no risk to ownership
(Theorem~\ref{thm:locked}, Corollary~\ref{cor:wait}), but no experiment here
isolates the profit of waiting: the ``cash locked half-suits immediately'' row is
inert by construction, and the locked-hold times of \S\ref{sec:locked} are
observational. \S\ref{app:human} states these three mechanisms in full and adds
the ask-ordering, declaration and endgame material, tagging each item by whether
it is a consequence of the rules, a measurement, or an untested suggestion.
